\section*{Chapter \ref{ch:b2:static-games} --- Static Games}

\begin{solution}{exo:b2:static-games:1}
% code: exercises/solutions.py iesds
For the column player, $C$ strictly dominates $R$ (2 against 1 against $T$, 1 against 0 against $B$): remove $R$. Then
$T$ strictly dominates $B$ for the row player (1 against 0 against both $L$ and $C$): remove $B$. Against $T$, the column
player prefers $C$ (2) to $L$ (0). The solution is $(T, C)$, with payoffs $(1, 2)$.
\end{solution}

\begin{solution}{exo:b2:static-games:2}
% code: exercises/solutions.py pure_nash
$(U, L)$: in column $L$ the row player's best is $U$ (3), and in row $U$ the column player's best is $L$ (3). $(M, C)$: 2 is
the row player's best in column $C$, 2 the column player's best in row $M$. $(D, R)$: 4 and 4. No other cell has both
properties: three pure Nash equilibria.
\end{solution}

\begin{solution}{exo:b2:static-games:3}
% figdata: contrib_summary
Strangers: from 18.5 to 57.5 per cent, 39 points. Fixed groups: from 37.5 to 85 per cent, 47.5 points. Strangers
without punishment in the final period, at 9.5 per cent, come closest to the prediction of zero.
\end{solution}

\begin{solution}{exo:b2:static-games:4}
Contributions that stay well above zero in the last period, when no future play can reward them, and among
players who will never meet again, refute it. Contributions that start positive but fall to zero by the last
period are consistent with players who learn the game, or who contribute early in the hope of encouraging
others, and the one-shot prediction applies to the final period only.
\end{solution}

\begin{solution}{exo:b2:static-games:5}
% code: exercises/solutions.py cournot
$q^* = 90/(n + 1)$ and $p^* = 100 - nq^*$. $n = 1$: 45 at 55, profit 2{,}025; $n = 2$: 30 each at 40, profit 900 each; $n = 4$: 18 each
at 28, profit 324; $n = 9$: 9 each at 19, profit 81.
\end{solution}

\begin{solution}{exo:b2:static-games:6}
% code: exercises/solutions.py public_good
$g^* = 12/4 = 3$ each, a total of 9. The efficient symmetric total is $3 \times 12/2 = 18$.
\end{solution}

\begin{solution}{exo:b2:static-games:7}
% code: exercises/solutions.py public_good
The ratio $2/(n + 1)$ falls: 0.67 for $n = 2$, 0.40 for $n = 4$, 0.18 for $n = 10$. Each person gives $12/(n + 1)$: 4, 2.4 and
1.09. A larger group raises the total only a little, from 8 to 9.6 to 10.9, while the efficient total grows in
proportion to $n$.
\end{solution}

\begin{solution}{exo:b2:static-games:8}
% code: exercises/solutions.py commons
$E^{\mathrm{N}} = 2 \times 80/3 = 53.33$; $E^* = 80/2 = 40$; over-use $53.33/40 - 1 = 33$ per cent. The fee is $80 \times 1/4 = 20$ per
boat-day: then $E = 2 \times 60/3 = 40$.
\end{solution}

\begin{solution}{exo:b2:static-games:9}
For each election, measure each candidate's position and the median voter's position on the same scale (from
surveys that place both). The proposition predicts that both candidates are close to the median and to each
other. Candidates who stay far apart in two-candidate races, consistently and not because of measurement
error, count against it, as does a gap that does not shrink when the median moves; before concluding, check
that the race is about one dimension and that candidates care only about winning.
\end{solution}

\begin{solution}{exo:b2:static-games:10}
(a) Suppose a \omterm{def:b2:static-games:nash-equilibrium}{Nash equilibrium} $s^*$ and let $s^*_i$ be the first of its strategies to be eliminated, at some
round, dominated by $t_i$ in the game that remains. Since no strategy of $s^*$ was removed before, $s^*_{-i}$ is
still present, so $u_i(t_i, s^*_{-i}) > u_i(s^*_i, s^*_{-i})$, which contradicts the equilibrium. (b) Player $i$
maximises $y - g_i + v(g_i + G_{-i})$; at an interior choice $v'(G) = 1$, and if $g_i = 0$ then $v'(G) \le 1$. With a positive
total some player contributes, so $v'(G^{\mathrm{N}}) = 1$. Total utility $ny - G + nv(G)$ is largest where $nv'(G^*) = 1$,
so $v'(G^*) = 1/n < 1 = v'(G^{\mathrm{N}})$; since $v'$ is decreasing, $G^* > G^{\mathrm{N}}$.
\end{solution}

\begin{solution}{exo:b2:static-games:11}
% code: exercises/solutions.py pure_nash
In matching pennies, at $(H, H)$ and $(T, T)$ the column player gains by switching; at $(H, T)$ and $(T, H)$ the row player
does. On $[0, 1]$, let $u_1 = -(s_1 - s_2)^2$ and $u_2 = (s_1 - s_2)^2$: player 1 wants to match, player 2 to be as far as
possible. If $s_1 \ne s_2$, player 1 gains by moving to $s_2$; if $s_1 = s_2$, player 2 gains by moving to 0 or 1. No profile is
an equilibrium.
\end{solution}

\begin{solution}{exo:b2:static-games:12}
If all three share a position $x$, each has $1/3$; one seller moving slightly to the longer side gains nearly $\max(x,
1 - x) \ge 1/2$. If exactly two share a position $x$ and the third is at $z > x$ (the case $z < x$ is symmetric), each of
the pair has $(x + z)/4$. Moving slightly left gives nearly $x$, moving slightly right nearly $(z - x)/2$; the first exceeds
$(x + z)/4$ when $z < 3x$ and the second when $z > 3x$. If $z = 3x$, the third seller gains by moving next to the pair,
from $1 - 2x$ to nearly $1 - x$. If all three are at different positions, the leftmost seller
gains by moving towards its neighbour: it keeps every buyer to its left and gains some between them. No profile
is an equilibrium.
\end{solution}

\begin{solution}{pb:b2:static-games:1}
% code: exercises/solutions.py commons quota_deviation
\textbf{Part I.}
1. $\pi_i = e_i(100 - e_i - E_{-i} - 20) = e_i(80 - e_i - E_{-i})$.
2. $80 - 2e_i - E_{-i} = 0$: $e_i = (80 - E_{-i})/2$.
3. No: the best effort falls as the others fish more.
4. The Cournot game, with effort for quantity.
5. One more boat-day lowers the value of every boat-day by one, and the boat bears that loss only on its own days.
6. Each boat would compare a fixed value with the cost, and its effort would not affect the others: no external
effect, and no over-fishing.
\textbf{Part II.}
7. $e = 80/6 = 13.33$ and $E = 66.67$.
8. $100 - 66.67 = 33.33$; profit $13.33 \times (33.33 - 20) = 177.78$ thousand kests.
9. $5 \times 177.78 = 888.89$.
10. The others' total is 53.33, and the \omterm{def:b2:static-games:best-response}{best response} is $(80 - 53.33)/2 = 13.33$.
11. $E = 10 \times 80/11 = 72.73$: more.
12. No: if every boat cut its effort to 8, each would earn $8 \times 40 = 320$, more than 177.78.
\textbf{Part III.}
13. Total profit $E(80 - E)$ is largest at $E^* = 40$.
14. Value 60 per boat-day; total profit $40 \times 40 = 1{,}600$.
15. $66.67/40 - 1 = 66.7$ per cent.
16. $1{,}600 - 888.89 = 711.11$ thousand kests a season.
17. Worse: each boat's catch would also lower future stocks, a further cost borne by all.
18. In equilibrium the five boats fish 66.7 per cent more than the efficient effort, and lose 711.11 thousand kests of
profit a season.
\textbf{Part IV.}
19. $f = 80 \times 4/10 = 32$: then $E = 5 \times 48/6 = 40$.
20. Each boat fishes 8 days and earns $8 \times (60 - 20 - 32) = 64$; the fee raises $32 \times 40 = 1{,}280$.
21. With the others at 32 days, the best effort is $(80 - 32)/2 = 24$, earning $24 \times 24 = 576$ against $8 \times 40 = 320$: a
gain of 256.
22. The expected fine must be at least 256: with detection probability one half, a fine of at least 512.
23. The fee needs an agency that knows the value and cost functions and the number of boats, and that sets the
fee in the boats' interest rather than its own; the boats keep only 64 each unless the revenue is returned.
The quota needs the boats to agree, a monitor, a fine that is paid, and a way to keep new boats out; the boats
keep the whole 1{,}600.
24. Over-fishing is 66.7 per cent of the efficient effort; a fee of 32 thousand kests per boat-day restores the
efficient effort of 40 boat-days.
\end{solution}
